\subsection{EXPONENTIAL MAPPINGS}

PROPOSITION 3. Let \(\mathbf{G}\) be a Lie group. There exists an exponential mapping \(\phi\) of \(\mathbf{G}\) with the following properties:

\begin{enumerate}
\def\labelenumi{(\roman{enumi})}
\item
  \(\phi\) is defined on an open subgroup U of the additive group L(G);
\item
  \(\phi (\mathbf{U})\) is an open subgroup of \(\mathbf{G}\) and \(\phi\) is an isomorphism of the analytic manifold \(\mathbf{U}\) onto the analytic manifold \(\phi (\mathbf{U})\) ;
\item
  \(\phi(nx) = \phi(x)^n\) for all \(x \in \mathbf{U}\) and all \(n \in \mathbf{Z}\) .
\end{enumerate}

Let \(\dot{\mathrm{L}}(\mathrm{G})\) be given a norm compatible with its topology and such that \(\| [x,y]\| \leqslant \| x\| \| y\|\) for \(x,y\) in \(\mathrm{L}(\mathrm{G})\) . Let \(\mathbf{G}_1\) be the Lie group defined by \(\mathrm{L}(\mathrm{G})\) . Let \(\psi = \mathrm{Id}_{\mathrm{G}_1}\) , which is an exponential mapping of \(\mathbf{G}_1\) . For all \(\mu >0\) , let \(\mathrm{L}_{\mu}\) be the set of \(x\in \mathrm{L}(\mathrm{G})\) such that \(\| x\| < \mu\) . Then, for \(\mu\) sufficiently small, \(\mathrm{L}_{\mu}\) is an open subgroup of the additive group \(\mathrm{L}(\mathrm{G}),\psi (\mathrm{L}_{\mu})\) is an open subgroup of \(\mathrm{G}_1\) (§ 4, no. 2, Lemma 3), \(\psi |\mathrm{L}_{\mu}\) is an isomorphism of analytic manifolds of \(\mathrm{L}_{\mu}\) onto \(\psi (\mathrm{L}_{\mu})\) and \(\psi (nx) = \psi (x)^n\) for all \(x\in \mathrm{L}_{\mu}\) and all \(n\in \mathbf{Z}\) . The \(\mathrm{L}_{\mu}\) form a fundamental system of neighbourhoods of 0 in \(\mathrm{L}(\mathrm{G})\) . By Theorem 1, there exist \(\mu\) and an open subgroup \(\mathbf{G}'\) of \(\mathbf{G}\) such that \(\psi (\mathrm{L}_{\mu})\) and \(\mathbf{G}'\) are isomorphic, whence the proposition.

PROPOSITION 4. Let G be a Lie group and \(\phi\) an injective exponential mapping of G. Suppose that \(p > 0\) . For all \(x, y\) in L(G),

\begin{enumerate}
\def\labelenumi{(\arabic{enumi})}
\tightlist
\item
\end{enumerate}

\[
x + y = \lim _ {n \rightarrow + \infty} p ^ {- n} \phi^ {- 1} (\phi (p ^ {n} x) \phi (p ^ {n} y))\tag{1}
\]

\[
[ x, y ] = \lim _ {n \rightarrow + \infty} p ^ {- 2 n} \phi^ {- 1} (\phi (p ^ {n} x) \phi (p ^ {n} y) \phi (- p ^ {n} x) \phi (- p ^ {n} y)).\tag{2}
\]

These are special cases of Proposition 4 of § 4, no. 3.

